\documentclass[11pt,letterpaper]{article}

\usepackage[margin=1in]{geometry}
\usepackage{amsmath,amssymb,amsthm}
\usepackage{physics}
\usepackage{booktabs}
\usepackage[numbers,sort&compress]{natbib}
\usepackage[hyperfootnotes=false]{hyperref}
\usepackage{cleveref}
\usepackage{float}
\usepackage{titlesec}
\usepackage{microtype}

% Section formatting
\titleformat{\section}{\large\bfseries\uppercase}{\thesection}{1em}{}
\titleformat{\subsection}{\normalsize\bfseries}{\thesubsection}{1em}{}

\hypersetup{
  colorlinks=true,
  linkcolor=black,
  citecolor=black,
  urlcolor=blue,
}

\newcommand{\hbarc}{\hbar c}
\newcommand{\Runiv}{R_{\mathrm{universe}}}
\newcommand{\Muniv}{M_{\mathrm{universe}}}
\newcommand{\Nuniv}{N_{\mathrm{universe}}}
\newcommand{\av}[1]{\left\langle #1 \right\rangle}

\title{An Entanglement Ansatz for the Mach--Sciama Relation:\\
Newton's Constant from Cosmic Boundary Conditions}
\author{}
\date{}

\begin{document}

\maketitle

\begin{abstract}
The relation $G \sim c^2 R/M$ between Newton's gravitational constant and
cosmic boundary conditions is well-known in the Mach--Sciama tradition: in
any flat critical-density cosmology, $G c^2 = (\text{const}) \cdot c^4 R/M$
follows from the Friedmann equation. What is \emph{not} explained in standard
cosmology is \emph{why} the universe should sit at this self-consistency point.
We propose a quantum entanglement mechanism --- orthogonal jerk creating
helical trajectories, with bond energy $\hbar c/r$ and torque geometry from
the Bisognano--Wichmann theorem --- that predicts the specific prefactor
$\gamma = 4/5$ from the modular Hamiltonian of the cosmic reduced density
matrix. The relation
\begin{equation}
  G \;=\; \frac{4}{5}\,\frac{c^2 \Runiv}{\Muniv}
\end{equation}
gives a value consistent with the measured $G$ to within current cosmic
parameter uncertainties.

\textbf{Scope.} This is not a free or local derivation of $G$. The
cosmic inputs $\Runiv$ and $\Muniv$ are themselves determined within
general-relativistic cosmology and contain $G$ implicitly
(Section~\ref{sec:scope}); applied to any local causal patch the
same formula misses by tens of orders of magnitude
(Section~\ref{sec:local_test}). The substantive content is therefore
not the numerical value of $G$ but \emph{(i)} the predicted prefactor
$\gamma = 4/5$, derived from the Racah-chain modular Hamiltonian at
low filling fraction and shown in
Section~\ref{sec:weight_forced} to be \emph{forced} by the regime
identification rather than fitted, and \emph{(ii)} the mechanism
itself, validated at the electron scale by the exact derivation of
the Bohr magneton in the companion paper. That companion result is
the load-bearing empirical leg of the program: it uses no $G$, no
fitted parameters, and no cosmic inputs.
\end{abstract}

\section*{The Empirical Anchor: A $G$-Free Test (Read First)}

The strongest evidence for the mechanism is \emph{not} the cosmic $G$ match
below --- which, as Section~\ref{sec:scope} explains, is a one-scale consistency
check entangled with $G$ through the cosmic inputs. The strongest evidence is an
\textbf{exact, parameter-free, $G$-free} result obtained by applying the
\emph{same} helical-jerk / modular-Hamiltonian mechanism at the electron Compton
scale: the electron's orthogonal-jerk trajectory is a current loop whose
magnetic moment is exactly the Bohr magneton,
\begin{equation}
  \mu = I A = \frac{e\,\omega_e}{2\pi}\,\pi\lambda_{C,e}^2 = \frac{e\hbar}{2 m_e}
      = \mu_B,
\end{equation}
using only $e$, $\hbar$, $m_e$, $c$, with $G$ nowhere present and no fitted
parameters or cosmic inputs (companion paper). The logic of this paper follows:
the mechanism is validated exactly and $G$-freely at the electron scale; what
follows is what the same mechanism implies for cosmic boundary conditions.
Gravity is an \emph{application} of a validated mechanism, not the primary
evidence for it. Read the cosmic number
(Section~\ref{sec:numerical}, below) in that light.

\section{Theoretical Foundation}

The ansatz rests on five independent results from quantum field theory
and conformal field theory.

\subsection{Bisognano--Wichmann Theorem (1975/1976)}
\label{sec:bw}

The modular Hamiltonian for a half-space in relativistic QFT is the generator
of Lorentz boosts perpendicular to the entangling surface:
\begin{equation}
  K = 2\pi \int_{x>0} \mathrm{d}^{d-1}x \; x \; T_{00}(x)
  \label{eq:bw}
\end{equation}
This is a \textbf{theory-independent} result. It applies to any relativistic QFT,
interacting or free. The modular Hamiltonian is always a weighted integral of
the stress-energy tensor $T_{00}$, with the weight being the distance $x$ from
the entangling surface.

\textbf{Physical meaning:} Entanglement creates a boost-rotation geometry, not a
time-translation. The nucleon's spin faces orthogonally to the displacement
toward the partner.

\subsection{Casini--Huerta--Myers (2011)}

For a spherical region of radius $R$ in a $d$-dimensional CFT, the modular
Hamiltonian is:
\begin{equation}
  K = 2\pi \int_{|\vec{x}|<R} \mathrm{d}^{d-1}x \; \beta(r) \; T_{00}(\vec{x})
  \label{eq:chm}
\end{equation}
where the weight function is:
\begin{equation}
  \beta(r) = \frac{R^2 - r^2}{2R}
  \label{eq:chm_weight}
\end{equation}
This is a parabola that vanishes at the boundary $r = R$ and is maximal at
the center $r = 0$. The result is derived by conformally mapping the ball to
a half-space and applying the BW theorem.

\subsection{Cardy--Tonni (2016)}

For an interval of length $l$ in 1+1D CFT, the entanglement Hamiltonian is
derived via conformal mapping from a punctured strip to an annulus:
\begin{equation}
  H_E = 2\pi \int_0^l \mathrm{d}x \; \beta(x) \; T_{00}(x)
  \label{eq:cardy_tonni}
\end{equation}
where the inverse temperature profile is:
\begin{equation}
  \beta(x) = \frac{l}{\pi}
  \frac{\cos(\pi x_0/l) - \cos(\pi x/l)}{\sin(\pi x_0/l)}
  \label{eq:cardy_tonni_beta}
\end{equation}
For the full interval ($x_0 = l/2$), $\cos(\pi/2) = 0$ and $\sin(\pi/2) = 1$, so:
\begin{equation}
  \beta(x) = -\frac{l}{\pi} \cos(\pi x/l)
  \label{eq:cardy_tonni_full}
\end{equation}
The energy scale is set by the inverse of the interval length: $E \sim \hbar c/l$.

\subsection{Bernard et al. (2024)}

The modular Hamiltonian for inhomogeneous free-fermion chains has a commuting
operator $T_A$ whose eigenvalues approximate the entanglement spectrum:
\begin{equation}
  \epsilon_k = \frac{2\pi \lambda_k}{\mathrm{scale} \cdot N^2}
  \label{eq:bernard_eigen}
\end{equation}
For the Racah chain (most general case), the full inverse temperature
$\beta(x) \propto x(L-x)$ is parabolic. When the filling fraction is low
($\rho \ll 1/2$), the Fermi velocity vanishes outside a narrow active region
$[x_-, L]$ near the boundary. Setting $u = (L-x)/L \in [0, 1-x_-/L]$ and
re-expressing $\beta$ on the active interval, the dominant term is
$\beta(r) \propto (1-r/R)^2$ to leading order in the ratio $(L-x_-)/L$
\citep{bernard2024}:
\begin{equation}
  \beta(r) = (1 - r/R)^2
  \label{eq:racah_weight}
\end{equation}
This $(1-r/R)^2$ form is established for \emph{specific integrable
free-fermion chains} as the leading term in the low-filling limit; it is not a
general theorem about arbitrary low-density states. Carrying it over to the
cosmic ensemble requires the separate (and unproven) identification discussed in
Section~\ref{sec:weight_forced} --- ``$\Omega_b \approx 0.049$ is small'' is
necessary but not sufficient justification.

\subsection{Huerta \& van der Velde (2023--2024)}

The modular Hamiltonian for a spherical region in $d$-dimensional CFT can be
obtained by dimensional reduction to the radial semi-infinite line. When a
$d$-dimensional free massless field is decomposed into angular modes, the
modular Hamiltonian decomposes as:
\begin{equation}
  K = \sum_{\ell, \vec{m}} K_{\ell, \vec{m}}
  \label{eq:huerta_decomp}
\end{equation}
where each mode contribution $K_{\ell, \vec{m}}$ is \textbf{exactly} the dimensional
reduction of the parent CFT modular Hamiltonian. The weight function $\beta(r)$
is \textbf{preserved} under dimensional reduction.

Critically, the 1D reduced theories are \textbf{non-conformal} (they carry a
$1/r^2$ potential from the angular Laplacian), yet their modular Hamiltonian
remains local in the energy density with the \textbf{same weight function} as the
$d$-dimensional CFT. This is proven for both scalar fields and Dirac fermions
\citep{huerta2023a, huerta2023b}.

\section{Circularity and Scope}
\label{sec:scope}

Before presenting the mechanism it is essential to be explicit about
what this calculation does and does not establish.

\subsection{The Mach--Sciama Relation Is Not Novel}

In a spatially flat universe at critical density, the Friedmann equation
gives $\rho_c = 3 H^2/(8\pi G)$. The total mass-energy in the Hubble
volume is then
\begin{equation}
  M_H \;=\; \rho_c \cdot \frac{4}{3}\pi R_H^3
       \;=\; \frac{c^2 R_H}{2 G},
       \qquad R_H = c/H,
\end{equation}
which immediately rearranges to $G = c^2 R_H/(2 M_H)$. A relation of the
form $G \sim c^2 R/M$ is therefore present in any standard cosmology;
it has been discussed since Sciama (1953) and is the structural premise
underlying Brans--Dicke theory and several modern emergent-gravity
programs. \emph{Reproducing this structural relation is not, by itself,
new physics.}

\subsection{The Inputs Are Not $G$-Independent}

In standard cosmology, both $\Runiv$ and $\Muniv$ are determined within
the GR/$\Lambda$CDM framework and depend on $G$:

\begin{itemize}
  \item $\Runiv$ is the proper distance to the particle horizon, computed
    from a Friedmann integral that contains $G$ through the matter
    density. Stellar age estimates (white dwarf cooling, globular cluster
    main-sequence turn-off) provide a partially $G$-independent age
    constraint, but the radius $\Runiv = c \cdot t_{\mathrm{age}}$ is at
    best a kinematic proxy.
  \item $\Muniv$ in $\Lambda$CDM is set by the closure condition
    $\Omega_{\mathrm{tot}} = 1$, which contains $G$ explicitly. Dark
    matter mass is determined by gravitational dynamics, so it cannot
    be measured independently of $G$. Only the baryonic component
    ($\Omega_b \approx 0.049$, from BBN and the CMB) is constrained
    by non-gravitational physics.
\end{itemize}

The numerical match between $\frac{4}{5} c^2 \Runiv/\Muniv$ and the
measured $G$ is therefore \emph{not} evidence that the entanglement
mechanism predicts the correct value of $G$ from $G$-free inputs. It
is evidence that the entanglement mechanism is consistent with the
Mach--Sciama relation to within the prefactor.

\subsection{Alternate Choices of $\Runiv$ and $\Muniv$}

One might hope to escape the circularity by picking inputs that are
less $G$-dependent. None of the obvious choices succeed:

\begin{table}[H]
  \centering
  \caption[Alternate choices of R and their G-independence.]{Alternate choices of $R$ and their $G$-independence.}
  \label{tab:R_choices}
  \begin{tabular}{p{4.4cm}p{3.4cm}p{6cm}}
    \toprule
    Choice for $R$ & $G$-independent? & What the claim becomes \\
    \midrule
    Particle horizon ($\sim$46 Gly) & No --- Friedmann integral with $G$ inside & Equivalent to the present paper \\
    Hubble radius $c/H_0$ & Partially --- $H_0$ from redshift-distance is the cleanest & Equivalent to "universe is at critical density" \\
    CMB last-scattering radius & No --- requires GR perturbation theory to interpret & Buys nothing \\
    Local causal diamond & Yes in principle & Promising; addressed in Sec.~\ref{sec:local_test} \\
    \bottomrule
  \end{tabular}
\end{table}

\begin{table}[H]
  \centering
  \caption[Alternate choices of M and their G-independence.]{Alternate choices of $M$ and their $G$-independence.}
  \label{tab:M_choices}
  \begin{tabular}{p{4.4cm}p{3.4cm}p{6cm}}
    \toprule
    Choice for $M$ & $G$-independent? & What the claim becomes \\
    \midrule
    Baryonic mass only & Yes --- CMB photon count $\times$ BBN $\eta$, no $G$ & Honest census $M_b \approx 1.5\times10^{53}$ kg; $G$ off by only $\sim 3\times$, not $\sim 20\times$ \\
    Baryon $+$ dark matter & No --- DM is gravitationally inferred & Reduces error but stays circular \\
    Total energy budget & No --- $\rho_c = 3H^2/(8\pi G)$ is explicitly circular & Identity, not derivation \\
    Komar/ADM/Misner-Sharp & No --- GR mass notions & Inside what we'd derive \\
    \bottomrule
  \end{tabular}
\end{table}

The most revealing entry is "baryonic mass only." If the relation
$G = (4/5)\, c^2\Runiv/\Muniv$ were a genuine prediction from
non-gravitational inputs, the value of $\Muniv$ that "works" should
be the one that comes from non-gravitational measurements (BBN,
atomic spectra, stellar populations). It isn't. The value that makes
the relation match $G$ requires including dark matter and dark
energy, both of which are gravitationally inferred. The match is
therefore even more circular than the headline calculation suggests:
most of the mass-energy "input" is itself a $G$-laden output.

\paragraph{A fully $G$-free evaluation.}
We can push harder and feed the formula \emph{only} inputs obtainable without
$G$: a radius from expansion kinematics, and a baryonic mass from CMB
photon-counting (the photon density $n_\gamma$ follows from the blackbody
temperature) times the BBN baryon-to-photon ratio $\eta$. An honest baryon
census in the particle-horizon volume gives $M_b \approx 1.5\times10^{53}$ kg ---
about $6\times$ larger than the $\Omega_b \times \Muniv$ figure, which silently
inherits the back-fit total. With $R$ the particle horizon, this purely $G$-free
pair yields $G_{\mathrm{pred}}/G_{\mathrm{meas}} \approx 3.2$ --- off by a factor
of $\sim 3$, \emph{not} the factor of $\sim 20$ that a naive $\Omega_b$ scaling
suggests. (With the Hubble radius instead, the same census is off by $\sim 33$,
since $G \propto 1/R^2$ at fixed density.)

A clean analytic check sharpens what the prefactor is doing. At the Hubble
radius and critical density, $c^2 R/M = 2G$ identically, so the $4/5$ prefactor
gives \emph{exactly}
\begin{equation}
  G_{\mathrm{pred}} = \tfrac{8}{5}\,G = 1.600\,G,
\end{equation}
overshooting the textbook Mach--Sciama coefficient $1/2$ by $8/5$. The
$0.21\%$ agreement of Section~\ref{sec:numerical} is recovered only by pairing
the particle horizon (about $3.2\times$ larger than the Hubble radius) with the
gravitationally inferred total mass. No purely $G$-free $(R, M)$ pair reproduces
the measured $G$; the closest is within a factor of $\sim 3$. All of these
evaluations are reproduced in \texttt{gfree\_inputs.py}.

\paragraph{Why $8/5$, and why it is not a circumference.}
It is tempting to read the overshoot as a hidden $2\pi$ --- a stray
circumference $2\pi r$ lurking in the geometry. It is not. The factorization is
exactly $8/5 = 2 \times (4/5)$, and the extra $2$ is the \emph{Schwarzschild}
factor, the same $2$ as in $R_H = 2GM/c^2$: at critical density the Hubble
radius \emph{is} the Schwarzschild radius of the mass it encloses. Equivalently
it is the $2$ that separates Einstein's coupling $8\pi G$ from the Newtonian
Poisson coefficient $4\pi G$, since $8\pi G = 2 \cdot 4\pi G$. The angular
factors cancel cleanly --- the $4\pi$ of the sphere volume against the $4\pi$
inside the Friedmann $8\pi$ --- so \emph{no $\pi$ survives} and there is no
$2\pi r$; this cancellation is verified symbolically in
\texttt{verify\_symbolic.py}. Notably, the mechanism's \emph{own} factor of
$2\pi$ --- the Bisognano--Wichmann/Unruh normalization $K = 2\pi\int x\,T_{00}$
(Section~\ref{sec:bw}) --- cancels out of the cosmic prefactor along with
$\hbar$, $m_p$, and $\kappa$ (the cancellation below
Eq.~\eqref{eq:final_g}); it therefore does \emph{not} source this Schwarzschild
$2$. The two would only meet through the local derivation of
Section~\ref{sec:local_derivation}: in the Jacobson construction the
gravitational $8\pi$ factorizes the \emph{other} way, as
$(2\pi)_{\mathrm{Unruh}} \times 4_{\mathrm{Bekenstein}}$, so reconciling the
mechanism's Unruh $2\pi$ with the gravitational coupling is precisely the
bond-density problem left open there. The $8/5$ is thus a clean diagnostic, not
a clue to a missing geometric factor.

\subsection{The Local Test, and What It Tells Us About the Relation}
\label{sec:local_test}

To understand what kind of statement $G = \gamma\, c^2 R/M$ actually
\emph{is}, apply it to a system that is not the entire observable
universe. Try Earth:
\begin{align}
  M_E &= 5.97 \times 10^{24}\;\mathrm{kg}, \\
  R_E &= 6.37 \times 10^{6}\;\mathrm{m}, \\
  G_{\mathrm{eff}}^{\mathrm{Earth}} &= \frac{4}{5} \frac{c^2 R_E}{M_E}
    \approx 7.7 \times 10^{25}\;\mathrm{m^3/(kg\,s^2)}.
\end{align}
This is $\sim 36$ orders of magnitude away from the measured
$G \approx 6.67 \times 10^{-11}$. Read naively, the formula appears
to ``fail'' the local test. But the more useful reading is: the
local test reveals what kind of relation this actually is. It is
\emph{not} a local law of the form ``$G$ depends on the local mass
and radius.'' It is a \emph{cosmic boundary-condition relation}.

\paragraph{What the local mismatch tells us.}
At the cosmic scale, $\Muniv/\Runiv \approx c^2/(2 G)$ because the
universe sits at critical density. The Mach--Sciama relation
$G \sim c^2 R/M$ is therefore an algebraic restatement of the
critical-density closure condition. It contains real physics --- the
predicted prefactor $4/5$ from the modular Hamiltonian --- and is
testable in the limited sense that the prefactor could have been
something else. But the relation as a whole only carries information
at the one scale where $M/R$ is forced by the closure condition. It
is not a local field equation.

\paragraph{What this implies for the program.}
Two things follow:
\begin{enumerate}
  \item The cosmic calculation should be read as a \emph{boundary
    condition statement}: the universe sits at critical density,
    and our mechanism predicts the prefactor in the resulting
    Mach--Sciama relation. This is real content. It is not, on its
    own, a local theory of gravity.
  \item A local theory of gravity from the same mechanism would
    require a separate derivation, in the spirit of
    Jacobson--Padmanabhan--Verlinde, from local causal horizons or
    entanglement thermodynamics. We sketch what that derivation
    would need in Section~\ref{sec:local_derivation}, and we do not
    yet have it.
\end{enumerate}

So the right framing is not ``the formula fails locally'' but ``the
formula is a cosmic boundary-condition statement, and a local
companion derivation is the open problem.'' The Bohr magneton
result (companion paper) is the closest thing we currently have to
a local empirical anchor for the underlying mechanism.

\subsection{What Is and Is Not Claimed}

\textbf{What is claimed:}
\begin{enumerate}
  \item There is a physically motivated entanglement mechanism
    (orthogonal jerk + helical trajectory + bond energy $\hbar c/r$)
    whose application at the electron Compton scale exactly
    reproduces the Bohr magneton, with no $G$ as input. This is the
    only clean numerical test of the mechanism (companion paper).
  \item The same mechanism, when applied to the cosmic ensemble with
    the standard $\Lambda$CDM inputs, reproduces the Mach--Sciama
    relation $G \sim c^2 R/M$ with a specific predicted prefactor
    $\gamma = 4/5$ derived from the modular Hamiltonian of free
    fermions on a Racah chain at low filling fraction.
\end{enumerate}

\textbf{What is not claimed:}
\begin{enumerate}
  \item This is \emph{not} a derivation of $G$ from $G$-free inputs.
    The cosmic parameters used are not $G$-independent, and the
    relation evaluated on any local patch (e.g.~Earth) gives
    nonsense by $\sim 36$ orders of magnitude.
  \item The $0.21\%$ point-estimate match between
    $(4/5) c^2\Runiv/\Muniv$ and the measured $G$ is therefore a
    consistency check at the one scale where the inputs are forced
    by the critical-density condition, not a precision prediction.
  \item The mechanism does not, on its own, replace general
    relativity, and the present treatment does not constitute a
    local field-theoretic derivation of Newtonian gravity from
    entanglement.
\end{enumerate}

The cleanest empirical leg of the program is therefore not this
gravity paper, but the Bohr magneton derivation at the electron
scale, which is local, uses no $G$, and matches an experimentally
measured quantity to high precision. The reader should weight the
program accordingly.

\section{The Bond Energy $E_{\mathrm{bond}} = \hbar c / r$}

\subsection{Dimensional Analysis}

The modular Hamiltonian $K$ is dimensionless (it appears in $\rho_A = e^{-K}/Z$).
The weight $\beta$ has dimensions $[L]$ (boost parameter = proper distance).
The energy density $T_{00}$ has dimensions $[E]/[L]^d$. The volume element
$\mathrm{d}^d x$ has dimensions $[L]^d$.

\begin{equation}
  K = 2\pi \int \mathrm{d}^d x \; \beta(x) \; T_{00}(x)
  \label{eq:dim_analysis}
\end{equation}

The integral has dimensions $[L]^d \cdot [L] \cdot [E]/[L]^d = [E] \cdot [L]$.
For $K$ to be dimensionless:
\begin{equation}
  \frac{[E] \cdot [L]}{\hbar c} = \mathrm{dimensionless}
  \implies E \sim \frac{\hbar c}{L}
  \label{eq:scaling}
\end{equation}
This is the universal scaling: the entanglement energy between two subsystems
at separation $L$ is $\hbar c/L$, up to a dimensionless coefficient.

\paragraph{Region entanglement vs.\ pairwise bonds (a genuine leap).}
The modular Hamiltonian describes the entanglement of a \emph{region} with its
\emph{complement}; it is not, on its face, a sum of pairwise bonds between
individual particles. Dimensional analysis fixes the scaling of the
region--complement entanglement energy at scale $L$, but the subsequent
treatment (Sections~5--6) decomposes the cosmic entanglement into pairwise terms
$C\,\kappa\hbar c/r$ summed over nucleon pairs. That pairwise decomposition is an
\emph{assumption}, not a theorem: it is the ansatz that the region-level modular
structure can be repackaged as additive two-body bonds with the same $\hbar c/r$
scaling. It is plausible for a dilute, weakly-correlated ensemble, but we do not
derive it from the modular Hamiltonian, and a referee is entitled to treat the
pairwise picture as a modeling choice rather than a consequence.

\subsection{Why Free Fermions Apply to Nucleons}

The Bernard et al. calculation works with free fermions. Nucleons interact via
the strong force. Why does the free fermion result apply?

\begin{enumerate}
  \item \textbf{BW theorem is theory-independent:} The modular Hamiltonian structure
    $K = 2\pi \int \beta T_{00}$ applies to any relativistic QFT, interacting
    or free. The specific form of $T_{00}$ does not change the local structure.

  \item \textbf{Dimensional analysis is theory-independent:} The scaling
    $E \sim \hbar c/L$ depends only on the dimensions of $T_{00}$ and $\beta$,
    which are universal.

  \item \textbf{Bernard et al. is a verification, not an assumption:} It confirms
    the universal structure in a concrete solvable model.

  \item \textbf{EFT argument (with a caveat):} The long-wavelength behavior is
    controlled by the effectively massless, low-energy sector, and the $1/r$
    scaling depends only on dimensions, not on the interactions. We caution,
    however, that integrating out QCD does \emph{not} turn nucleons into free
    fermions: they remain massive composite particles with residual
    interactions. The free-fermion modular Hamiltonian is therefore used as a
    tractable \emph{model} of the low-filling sector --- justified by points
    1--2, where the structure and scaling are theory-independent --- not by any
    claim that nucleons are literally non-interacting.

  \item \textbf{Cross-check:} At the nucleon scale ($r \sim 1$ fm),
    $\hbar c/r \sim 200$ MeV, matching $\Lambda_{\mathrm{QCD}}$. This agreement at
    both nucleon and cosmic scales is independent confirmation.
\end{enumerate}

\subsection{Numerical Verification}

\begin{table}[H]
  \centering
  \caption[Bond energy across scales]{Bond energy $E_{\mathrm{bond}} = \hbar c/r$ across scales.}
  \label{tab:bond_energy}
  \begin{tabular}{llll}
    \toprule
    Scale & $r$ & $E_{\mathrm{bond}} = \hbar c/r$ & Physical interpretation \\
    \midrule
    Nucleon & 1 fm & $3.2 \times 10^{-11}$ J (200 MeV) & Strong force scale ($\Lambda_{\mathrm{QCD}}$) \\
    Earth & 6371 km & $5.0 \times 10^{-33}$ J & $3.1 \times 10^{-14}$ eV \\
    Cosmic & $4.4 \times 10^{26}$ m & $7.2 \times 10^{-53}$ J & Gravitational binding scale \\
    \bottomrule
  \end{tabular}
\end{table}

\section{The Geometric Factor $\gamma = 4/5$}

\subsection{The 1D Racah Chain Result}

The cosmic filling fraction is $\rho = \Omega_b = 0.049$ (the baryon fraction).
In the Racah chain framework, this maps to parameters:
\begin{equation}
  x_1 = \frac{\Omega_{\mathrm{DM}}}{\Omega_b} = 5.47, \quad
  x_2 = \frac{\Omega_{\mathrm{DE}}}{\Omega_b} = 13.94, \quad
  x_3 = \frac{\Omega_{\mathrm{DM}}}{\Omega_{\mathrm{DE}}} = 0.39
  \label{eq:racah_params}
\end{equation}
The Fermi velocity $v_F(x)$ is non-zero only on the active region
$[x_-, x_+] = [0.868, 1.000]$, a narrow band near the boundary.
The modular Hamiltonian weight on this region is quadratic:
\begin{equation}
  \beta_{\mathrm{1D}}(x) = (1 - x/L)^2
  \label{eq:racah_1d}
\end{equation}

\subsection{The 3+1D Extension}

The Huerta \& van der Velde theorem proves that the 1D modular Hamiltonian
on the radial line is the dimensional reduction of the $d$-dimensional sphere.
The weight function $\beta(r)$ is \textbf{preserved} under dimensional reduction.

In 3+1D, the spherical volume element $r^2 \mathrm{d}r$ modifies the weighted average.
The weighted average of $1/r$ with the quadratic weight is:
\begin{equation}
  \av{\frac{1}{r}}_{\mathrm{3D}}
  = \frac{\int_0^R \beta(r) \cdot r \, \mathrm{d}r}
         {\int_0^R \beta(r) \cdot r^2 \, \mathrm{d}r}
  \label{eq:weighted_avg}
\end{equation}

For $\beta(r) = (1 - r/R)^2$:

\noindent\textbf{Numerator:}
\begin{equation}
  \int_0^R (1 - r/R)^2 \cdot r \, \mathrm{d}r
  = R^2 \int_0^1 u^2(1-u) \, \mathrm{d}u
  = R^2 \left(\frac{1}{3} - \frac{1}{4}\right)
  = \frac{R^2}{12}
  \label{eq:numerator}
\end{equation}

\noindent\textbf{Denominator:}
\begin{equation}
  \int_0^R (1 - r/R)^2 \cdot r^2 \, \mathrm{d}r
  = R^3 \int_0^1 u^2(1-u)^2 \, \mathrm{d}u
  = R^3 \left(\frac{1}{3} - \frac{1}{2} + \frac{1}{5}\right)
  = \frac{R^3}{30}
  \label{eq:denominator}
\end{equation}

\noindent\textbf{Result:}
\begin{equation}
  \av{\frac{1}{r}} = \frac{R^2/12}{R^3/30} = \frac{5}{2R}
  \label{eq:avg_result}
\end{equation}

The geometric factor is:
\begin{equation}
  \gamma = \frac{2}{\av{1/r} \cdot R} = \frac{2}{5/2} = \frac{4}{5}
  \label{eq:gamma}
\end{equation}

\subsection{Sensitivity to the Weight Exponent}

The quadratic weight $a = 2$ is determined by the Racah chain physics
(low filling fraction $\rho = 0.049$), not by fitting to $G$. A scan
of nearby exponents shows the result is sensitive: the error grows
rapidly as $a$ deviates from 2.

\begin{table}[H]
  \centering
  \caption[Sensitivity to weight exponent]{Sensitivity of $\gamma$ to the weight exponent $a$ in $\beta(r) = (1-r/R)^a$.}
  \label{tab:sensitivity}
  \begin{tabular}{cccc}
    \toprule
    $a$ & $\av{1/r}\,R$ & $\gamma$ & Error vs measured $G$ \\
    \midrule
    1.8 & 2.400 & 0.8333 & $4.39\%$ \\
    1.9 & 2.450 & 0.8163 & $2.26\%$ \\
    \textbf{2.0} & \textbf{2.500} & \textbf{0.8000} & \textbf{$0.21\%$} \\
    2.1 & 2.550 & 0.7843 & $1.75\%$ \\
    2.2 & 2.600 & 0.7692 & $3.64\%$ \\
    \bottomrule
  \end{tabular}
\end{table}

The continuous family $\beta(r) = (1-r/R)^a$ admits the closed form
\begin{equation}
  \gamma(a) = \frac{2}{\av{1/r}\,R} = \frac{4}{a+3},
  \label{eq:gamma_of_a}
\end{equation}
so $a = 2 \Rightarrow \gamma = 4/5$, and the scan above is sensitive near
$a = 2$.

\paragraph{Moment decomposition of $\gamma$.}
The weighted average $\av{1/r}$ admits an exact decomposition in terms of the
first moment $\av{r}_\beta = R/(a+2)$ (the mean position of the bare weight
$\beta$):
\begin{equation}
  \av{\frac{1}{r}} = \frac{1}{2\,\av{r}_\beta} + \frac{1}{2R}.
  \label{eq:moment_decomp}
\end{equation}
Substituting into $\gamma = 2/(\av{1/r}\,R)$ gives
\begin{equation}
  \boxed{\;\gamma = \frac{4}{R/\av{r}_\beta + 1}\;.}
  \label{eq:gamma_moment}
\end{equation}
This is the physical content of the prefactor: $\gamma$ is determined by
how deep inside the volume the mean entanglement weight sits. For the
Racah weight ($a = 2$), the mean is at $\av{r}_\beta = R/4$ --- the inner
quarter of the radius --- giving $R/\av{r}_\beta = 4$ and $\gamma = 4/5$.

This decomposition makes explicit how each moment of the
$\mathrm{SL}(2,\mathbb{R})$ structure enters the final result:
\begin{itemize}
  \item \textbf{0th moment} (normalization): absorbed into the geometric
    coefficient $\kappa$, which cancels in the final expression for $G$
    (Section~\ref{sec:g_emerges}).
  \item \textbf{1st moment} (mean position $\av{r}_\beta$): determines
    $\gamma$ through $R/\av{r}_\beta$. This is the only moment that
    survives into the prefactor.
  \item \textbf{2nd moment} (variance): a slave variable of $a$
    (Section~\ref{sec:moment_hierarchy}), which is already set by
    $\av{r}_\beta$. It carries no independent information.
\end{itemize}

The continuous scan is, however, only illustrative: the physics does
not admit an arbitrary exponent. The admissible weights are the two
\emph{discrete} universal CFT results --- the vacuum (CHM) weight and the
low-filling (Racah) weight --- which give $16/15$ and $4/5$ respectively. The
real question is therefore not ``what value of $a$ fits $G$?'' but ``which
discrete regime describes the cosmic ensemble?'' --- whether $a = 2$ is
\emph{forced} by physics or merely \emph{found} to fit. That is the subject of
the next subsection. But first we note a structural reason why higher-order
corrections to the weight cannot arise.

\subsection{Why the Moment Hierarchy Terminates at Rank 2}
\label{sec:moment_hierarchy}

The weight $\beta(r)$ determines $\gamma$ through the weighted average
$\av{1/r}$. One might worry that the quadratic form $\beta = (1-r/R)^2$ is only
the leading term, and that higher-order corrections --- a cubic
$(1-r/R)^3$ admixture, say --- could shift $\gamma$ away from $4/5$.

No such correction exists, for a reason that is algebraic rather than
numerical.

\paragraph{The $\mathrm{SL}(2,\mathbb{R})$ constraint.}
The modular Hamiltonian is constructed from the generators of the
residual conformal symmetry of the radial theory. Huerta \& van der Velde
show that dimensional reduction from $d$ dimensions to the radial
semi-infinite line preserves an $\mathrm{SL}(2,\mathbb{R})$ subalgebra of the
parent conformal group $\mathrm{SO}(d,2)$. The modular flow generator is a
Noether charge of this residual symmetry. $\mathrm{SL}(2,\mathbb{R})$ has
exactly three generators:
\begin{enumerate}
  \item the \emph{dilation} generator $D$ (rank~0: sets the overall scale),
  \item the \emph{translation} generator $P$ (rank~1: sets the center/dipole),
  \item the \emph{special conformal} generator $K$ (rank~2: sets the
    quadratic curvature of the weight).
\end{enumerate}
There is no fourth generator. A rank-3 invariant (which would source a
cubic correction to $\beta$) does not exist in $\mathrm{SL}(2,\mathbb{R})$.

\paragraph{Derivation: higher moments are slave variables.}
Consider the family $\beta(r) = (1 - r/R)^a$ on $[0, R]$.
The $k$th moment of this weight distribution (normalized to a
probability measure) is
\begin{equation}
  \mu_k(a) \;=\;
  \frac{\int_0^R r^k\,(1-r/R)^a\,\mathrm{d}r}
       {\int_0^R (1-r/R)^a\,\mathrm{d}r}
  \;=\; R^k\,\frac{B(k+1,\,a+1)}{B(1,\,a+1)}
  \;=\; R^k\,\frac{k!\;\Gamma(a+1)}{\Gamma(k+a+2)}\,(a+1),
  \label{eq:moments_beta}
\end{equation}
where $B$ is the Euler beta function. For fixed $a$ (which the
$\mathrm{SL}(2,\mathbb{R})$ structure determines), every moment $\mu_k$
is a fixed rational function of $a$ and $R$ --- there is no free parameter
beyond these two. Explicitly:
\begin{align}
  \mu_0 &= 1 \quad\text{(normalization)}, \label{eq:moment0} \\
  \mu_1/R &= \frac{1}{a+2} \quad\text{(mean)}, \label{eq:moment1} \\
  \mu_2/R^2 &= \frac{2}{(a+2)(a+3)} \quad\text{(second moment)}, \label{eq:moment2} \\
  \mu_3/R^3 &= \frac{6}{(a+2)(a+3)(a+4)} \quad\text{(third moment = skewness source)}, \label{eq:moment3} \\
  \mu_4/R^4 &= \frac{24}{(a+2)(a+3)(a+4)(a+5)} \quad\text{(fourth moment = kurtosis source)}. \label{eq:moment4}
\end{align}
The pattern is immediate: $\mu_k/R^k = k!\,/\prod_{j=2}^{k+1}(a+j)$.
Once $a$ is fixed (by the rank-2 generator), moments 3, 4, and all
higher moments are \emph{determined}. They are not independent degrees
of freedom; they are slave variables of the exponent $a$.

\paragraph{Consequence: an exhaustive dichotomy.}
Since $\gamma(a) = 4/(a+3)$ (\cref{eq:gamma_of_a}), and the
$\mathrm{SL}(2,\mathbb{R})$ structure admits only two physically
distinct regimes --- vacuum ($a = 1$, the CHM weight) and low-filling
($a = 2$, the Racah weight) --- the prefactor takes exactly one of
two values:
\begin{equation}
  \boxed{\;\gamma \;\in\;
    \left\{\;\frac{16}{15}\;\text{(vacuum)},\quad
            \frac{4}{5}\;\text{(low filling)}\;\right\}.}
  \label{eq:dichotomy}
\end{equation}
The ratio between them is $\gamma_{\mathrm{CHM}}/\gamma_{\mathrm{Racah}}
= 4/3$, i.e.\ a $33\%$ split. No intermediate value is admissible within
the conformal family; no higher-order correction can move $\gamma$ within
or beyond this pair.

The consequence for $G$ is equally sharp. With $G = \gamma\,c^2\Runiv/\Muniv$,
the mechanism predicts one of exactly two values:
\begin{equation}
  G \;\in\; \left\{\;
    \frac{16}{15}\,\frac{c^2\Runiv}{\Muniv},\quad
    \frac{4}{5}\,\frac{c^2\Runiv}{\Muniv}
  \;\right\}.
  \label{eq:G_dichotomy}
\end{equation}
The measured $G$ agrees with the low-filling branch to $0.21\%$ and
disagrees with the vacuum branch by $33\%$. This is the paper's
sharpest falsifiable claim: the cosmic ensemble must sit in the
low-filling regime, not the vacuum, and the moment hierarchy admits
no third option.

\paragraph{Why there is no rank 3.}
The termination is not an approximation. The $\mathrm{SL}(2,\mathbb{R})$
algebra satisfies $[D, P] = P$, $[D, K] = -K$, $[P, K] = 2D$.
A putative rank-3 generator $T$ would need to satisfy
$[D, T] = \lambda T$ for some weight $\lambda$, but the Casimir structure
of $\mathrm{SL}(2,\mathbb{R})$ (rank-1 Lie algebra, one independent
Casimir $C_2 = D^2 - \tfrac{1}{2}(PK + KP)$) cannot accommodate
a fourth independent generator. The representation theory is three-dimensional;
any attempt to add a fourth generator either decomposes into the existing
three or enlarges the algebra to $\mathrm{SL}(3,\mathbb{R})$ or beyond, which
would require additional spatial dimensions or internal degrees of freedom
not present in the radial reduction.

\paragraph{Note on symmetry breaking.} The conformal symmetry \emph{is}
broken by the finite filling fraction $\rho = 0.049$. But the breaking enters
through the Fermi velocity profile (Section~\ref{sec:weight_forced}), which
modifies the quadratic coefficient --- i.e., it changes the value of the
rank-2 generator's contribution, not its rank. The filling fraction shifts $a$
within the space $\{1, 2\}$ (CHM vs.\ Racah); it does not open a rank-3
channel.

\subsection{Was the Quadratic Weight Forced or Found?}
\label{sec:weight_forced}

The fairest objection to this paper is: ``you wrote down a quadratic
weight, the algebra gave $4/5$, and that matches measured $G$. How
do we know you did not select the weight to make this work?'' The
answer requires looking at what universal CFT results actually say
and where they lead.

\paragraph{Two universal results, two weights.}
There are two genuinely theory-independent CFT results about modular
Hamiltonians on spherical regions.

\begin{enumerate}
  \item \textbf{Casini--Huerta--Myers (CHM, vacuum case).} For the
    vacuum reduced to a ball of radius $R$ in any $d$-dimensional
    CFT, the modular Hamiltonian is the boost generator with weight
    \begin{equation}
      \beta_{\mathrm{CHM}}(r) = \frac{R^2 - r^2}{2R}.
      \label{eq:chm_weight_forced}
    \end{equation}
    This vanishes \emph{linearly} as $r \to R$
    ($\beta \sim (R-r)$), is universal across CFTs, and applies to
    the vacuum state.
  \item \textbf{Bernard et al. (low-filling fermion case).} For free
    fermions on an inhomogeneous chain at filling fraction
    $\rho \ll 1/2$, the modular Hamiltonian's commuting operator
    has Fermi velocity vanishing on a narrow active region, and the
    leading weight is \emph{quadratic}:
    \begin{equation}
      \beta_{\mathrm{Racah}}(r) = (1 - r/R)^2.
      \label{eq:racah_weight_forced}
    \end{equation}
    This vanishes \emph{quadratically} as $r \to R$ and applies to
    a finite-density fermionic ground state, not to the vacuum.
\end{enumerate}

\paragraph{What the two weights predict.}
If we plug $\beta_{\mathrm{CHM}}$ into the same 3D weighted-average
machinery (\cref{eq:weighted_avg}):
\begin{align}
  \int_0^R \beta_{\mathrm{CHM}}(r) \cdot r \,dr
   &= \tfrac{R}{2}\!\!\int_0^R r\,dr - \tfrac{1}{2R}\!\!\int_0^R r^3\,dr
    = \tfrac{R^3}{4} - \tfrac{R^3}{8} = \tfrac{R^3}{8}, \\
  \int_0^R \beta_{\mathrm{CHM}}(r) \cdot r^2 \,dr
   &= \tfrac{R}{2}\!\!\int_0^R r^2\,dr - \tfrac{1}{2R}\!\!\int_0^R r^4\,dr
    = \tfrac{R^4}{6} - \tfrac{R^4}{10} = \tfrac{R^4}{15}, \\
  \av{1/r}_{\mathrm{CHM}} &= \tfrac{15}{8R}, \quad
  \gamma_{\mathrm{CHM}} = \tfrac{2}{15/8} = \tfrac{16}{15} \approx 1.067.
\end{align}
The vacuum (CHM) weight predicts $\gamma_{\mathrm{CHM}} = 16/15$,
which is \emph{wrong} by $33\%$. The low-filling (Racah) weight
predicts $\gamma_{\mathrm{Racah}} = 4/5$, which agrees with measured
$G$ to $0.21\%$.

\paragraph{Why the Racah weight is the relevant one.}
The two weights apply to different physical situations. CHM is the
vacuum modular Hamiltonian; Racah is the modular Hamiltonian of a
finite-density fermionic state at low filling. The cosmic ensemble
is not the vacuum. It contains baryons at filling fraction
$\Omega_b \approx 0.049$, far below $1/2$. This places it
unambiguously in the Bernard regime, not the CHM regime. \emph{We
are forced to use $\beta_{\mathrm{Racah}}$, not because $4/5$ is
the answer we wanted, but because the cosmic ensemble is a
low-filling state, not a vacuum.}

\paragraph{What is forced and what is not.}
The honest accounting:

\begin{itemize}
  \item \textbf{Forced.} If the cosmic ensemble is correctly
    described as a low-filling free-fermion ground state in the
    Racah-chain framework, then the quadratic weight
    $\beta(r) = (1-r/R)^2$ follows from Bernard et al., and the
    weighted-average machinery forces $\gamma = 4/5$ by simple
    integration. There is no remaining freedom: the
    $\mathrm{SL}(2,\mathbb{R})$ structure
    (Section~\ref{sec:moment_hierarchy}) forbids higher-order
    corrections to the weight, so neither skewness nor kurtosis of
    the weight profile can be adjusted independently.
  \item \textbf{Not forced.} The identification of the cosmic
    ensemble with a Racah-chain low-filling state. This is the
    fragile bridge, and it decomposes into three separate, individually
    unproven leaps:
    \begin{enumerate}
      \item \emph{Why fermionic?} The cosmic reduced density matrix
        contains bosons (photons, gluons, the Higgs, gravitons) as
        well as fermions. We have no argument for why its modular
        Hamiltonian should be approximated by a free-\emph{fermion}
        model rather than a mixed or bosonic one.
      \item \emph{Why $\Omega_b$ sets the filling?} We map the filling
        fraction to the baryon fraction $\Omega_b \approx 0.049$, but
        the matter fraction $\Omega_m \approx 0.31$ (also below $1/2$)
        or some other parameter would be at least as defensible. The
        specific mapping to the Racah parameters $x_1, x_2, x_3$ in
        \cref{eq:racah_params} is not derived.
      \item \emph{Why the Racah chain?} The Racah chain is the most
        general \emph{integrable} inhomogeneous free-fermion chain.
        The universe is not an integrable spin chain, and we do not
        argue why integrability is a good approximation at cosmic
        scales, only that it furnishes a tractable solvable model with
        the right qualitative low-filling behavior.
    \end{enumerate}
    Any one of these can be rejected, and with it the prefactor.
  \item \textbf{Therefore.} The claim is: \emph{conditional on the
    Racah-chain low-filling identification, the prefactor $4/5$ is
    forced.} The conditional is the load-bearing physics, not the
    arithmetic.
\end{itemize}

\paragraph{What would falsify the prefactor.}
The framework makes the prefactor a function of the modular weight,
which is a function of the regime. A different weight (linear or
hybrid) would give a different prefactor; but a \emph{cubic}
correction cannot arise within the $\mathrm{SL}(2,\mathbb{R})$ family
(Section~\ref{sec:moment_hierarchy}), so the only admissible alternative
is the CHM weight ($\gamma = 16/15$), which is $33\%$ wrong.
Reproducing the cosmic Mach--Sciama relation with a different prefactor
would falsify the specific claim of this paper. Reproducing it with no
identifiable weight at all --- pure $\gamma = 1$ --- would falsify
the entire mechanism.

\section{The Mechanism}

\subsection{Cosmic Entanglement Structure}

In the cosmic reduced density matrix, every nucleon is entangled with every
other nucleon. The angular momentum budget per nucleon is:
\begin{equation}
  s = \frac{\sqrt{3}}{2} \hbar
  \label{eq:spin_budget}
\end{equation}
This budget is distributed across $\Nuniv$ partners, with the
allocation following the modular Hamiltonian structure ($\hbar c/r$).

\subsection{Torque Geometry}

The BW theorem shows that the modular flow is a boost rotation, not a time
translation. The nucleon's spin faces orthogonally to the displacement toward
the partner. This is the ``torque geometry'' of entanglement.

\subsection{Unresolved Entanglement}

When the per-pair coupling between two particles is far below unity, the pair
has not collapsed to a definite correlation. The entanglement amplitude still
exists but remains unresolved. For gravity, this coupling is
$\alpha_{\mathrm{grav}} \approx 5.9 \times 10^{-39}$, so essentially all
nucleon pairs are in this unresolved regime.

This is not zero force. The unresolved amplitude carries the cosmic
entanglement field: each nucleon is entangled with $N_{\mathrm{universe}}$
other nucleons. The balance of angular momentum is the rest of the universe.
In QFT language, this is the reduced density matrix obtained by tracing out
the complement. The universe's isotropic entanglement field is the baseline;
any local mass breaks the isotropy.

\subsection{Gravity as Anisotropy}

The boost rotation back toward the mass is the gravitational force. Gravity
is not a fundamental force---it is the anisotropy in the cosmic entanglement
field created by local mass concentrations.

\subsection{No Angular Average}
\label{sec:no_angular_avg}

A conventional field theory would compute the gravitational force as an
angular average: integrate the force density over solid angle, project onto
the radial direction via $\cos\theta$, and obtain a statistical expectation
value. We do not do this, and the reason is not laziness but principle.

Each nucleon pair is a \emph{definite} entanglement bond with a definite
separation vector $\mathbf{r}_{ij}$. The bond carries a definite force
$C\,\kappa\,\hbar c / r_{ij}^2$ along $\hat{\mathbf{r}}_{ij}$. There is
no ensemble of orientations to average over: the pair either exists or it
does not. The total force on a nucleon is the \emph{sum} over all its
partners, not an \emph{average} over angles:
\begin{equation}
  \mathbf{F}_i = \sum_{j \neq i} C\,\kappa\,\frac{\hbar c}{r_{ij}^2}
                 \;\hat{\mathbf{r}}_{ij}.
  \label{eq:force_sum}
\end{equation}
The isotropic part of this sum cancels by symmetry (the universe is
isotropic). The residual is the force toward the local mass
concentration. No $\cos\theta$ projection is performed; no angular
measure is introduced.

This distinction matters because an angular average is ontologically
violent: it replaces definite bonds with a statistical smear, destroying
the pairwise structure that the modular Hamiltonian provides. In this
framework, the entanglement bond between nucleon $i$ and nucleon $j$ is
a real object --- it carries a definite coupling $C$, a definite energy
$C\kappa\hbar c/r_{ij}$, and a definite direction $\hat{\mathbf{r}}_{ij}$.
The force is the derivative of a sum of definite terms, not the
projection of a continuous field. The pair counting
$(M/m_p)(m/m_p)$ in the ansatz (Section~\ref{sec:g_emerges}) is exact
for this reason: it counts bonds, not field modes.

\section{The Ansatz}

\textbf{Terminology note.} Throughout this section we use the symbol $C$ to
denote the per-pair \emph{coupling fraction} (the strength with which two
nucleons share entanglement). This is \emph{not} the Wootters concurrence of
two-qubit quantum information theory, and the two should not be conflated; an
earlier draft used ``concurrence'' for $C$, which was misleading.

\subsection{Setup}

Each nucleon carries angular momentum $s = \sqrt{3}/2 \cdot \hbar$. From
the modular Hamiltonian (Sections~1--2), the bond energy between two
nucleons at separation $r$ is:
\begin{equation}
  E_{\mathrm{bond}} = \kappa \frac{\hbar c}{r}
  \label{eq:bond_energy}
\end{equation}
where $\kappa$ is an order-unity geometric coefficient. The BW theorem
sets the torque geometry: the nucleon faces orthogonally to the
displacement toward the partner. The isotropic cosmic background
cancels, leaving only the anisotropy from local mass concentrations.

\subsection{Per-Pair Coupling from Energy Balance}

The per-pair coupling $C$ is determined by the total entanglement
energy of the universe. The universe's mass-energy is stored as
entanglement energy across all nucleon pairs:
\begin{equation}
  E_{\mathrm{entanglement}}
  = \frac{\Nuniv^2}{2} \cdot C \cdot \kappa \hbar c \cdot \av{\frac{1}{r}}
  \label{eq:total_entangle}
\end{equation}
where $\av{1/r} = 5/(2\Runiv)$ from the quadratic weight.
Equating to the universe's mass-energy $\Muniv c^2$:
\begin{equation}
  \frac{N^2}{2} \cdot C \cdot \kappa \hbar c \cdot \frac{5}{2R} = M c^2
  \label{eq:energy_balance}
\end{equation}

\paragraph{This is a constraint, not a dynamical law.}
\Cref{eq:energy_balance} is the weakest step in the chain, and we flag it as
such. It is an \emph{imposed} equality --- ``the universe's mass-energy equals
its total pairwise entanglement energy'' --- not a consequence derived from a
dynamical principle. It is in the spirit of holographic equipartition
\citep{padmanabhan2010}, where the bulk energy is matched to a surface degree-of-freedom
count, but we do not derive it from equipartition or from any equation of
motion. A reader who does not grant this matching condition does not get $C$,
and hence does not get the force law that follows. Supplying a dynamical
justification --- ideally the same one that would fix the local bond density
$\eta$ in Section~\ref{sec:local_derivation} --- is part of the open problem.

Solving for $C$ and substituting $N = M/m_p$:
\begin{equation}
  C = \frac{4}{5} \frac{m_p^2 c \Runiv}{\kappa \Muniv \hbar}
  \label{eq:coupling_frac}
\end{equation}

\subsection{$G$ Emerges}
\label{sec:g_emerges}

The force per nucleon pair is the derivative of the bond energy:
\begin{equation}
  F_{\mathrm{pair}} = -\frac{\mathrm{d}}{\mathrm{d}r}
    \left(C \cdot \kappa \frac{\hbar c}{r}\right)
  = C \cdot \kappa \frac{\hbar c}{r^2}.
  \label{eq:pair_force}
\end{equation}
This is the force along the definite bond direction $\hat{\mathbf{r}}_{ij}$
(Section~\ref{sec:no_angular_avg}), not a projection or angular average.
The total force between two masses $M$ and $m$ at separation $r$ is the sum
over all $N_M \times N_m$ definite bonds:
\begin{equation}
  F = \frac{M}{m_p} \cdot \frac{m}{m_p} \cdot C \cdot \kappa \frac{\hbar c}{r^2}
  \label{eq:total_force}
\end{equation}
Substituting \cref{eq:coupling_frac}:
\begin{equation}
  F = \frac{4}{5} \cdot M \cdot m \cdot \frac{c^2 \Runiv}{\Muniv} \cdot \frac{1}{r^2}
  \label{eq:newton_form}
\end{equation}
Comparing to Newton's law $F = G M m / r^2$:
\begin{equation}
  G = \frac{4}{5} \frac{c^2 \Runiv}{\Muniv}
  \label{eq:final_g}
\end{equation}
\textbf{Note:} $\hbar$, $m_p$, and $\kappa$ all cancel. The final expression depends
only on $c$, $\Runiv$, and $\Muniv$.

\section{Numerical Consistency Check}
\label{sec:numerical}

Using cosmic inputs:
\begin{align}
  c &= 2.998 \times 10^8 \; \mathrm{m/s} \\
  \Runiv &= 4.4 \times 10^{26} \; \mathrm{m} \\
  \Muniv &= 4.73 \times 10^{53} \; \mathrm{kg}
\end{align}
the predicted relation gives
\begin{equation}
  G_{\mathrm{predicted}}
  = \frac{4}{5} \frac{(2.998 \times 10^8)^2 \cdot (4.4 \times 10^{26})}
                     {4.73 \times 10^{53}}
  = 6.688 \times 10^{-11} \; \mathrm{m^3/(kg{\cdot}s^2)},
  \label{eq:g_derived}
\end{equation}
to be compared with $G_{\mathrm{measured}} = 6.674 \times 10^{-11}\;
\mathrm{m^3/(kg{\cdot}s^2)}$. The ratio is $1.0021$.

This check should be read in the light of Section~\ref{sec:scope}. The central
value carries two very different uncertainties, propagated through
$\delta G/G = \sqrt{(\delta R/R)^2 + (\delta M/M)^2}$:

\begin{itemize}
  \item \textbf{Measurement uncertainty} (definitions of $R$, $M$ fixed). With
    the particle-horizon radius and the critical-density mass at Planck-level
    precision, $\delta R/R \sim 2\%$ and $\delta M/M \sim 3\%$ give
    $\delta G/G \sim 4\%$, i.e.\ $G_{\mathrm{predicted}} = (6.7 \pm 0.3)\times
    10^{-11}$. The measured value sits comfortably inside this band --- but the
    band is $\sim 20\times$ wider than the $0.21\%$ central offset, so the
    agreement is a \emph{consistency}, not a precision test.
  \item \textbf{Definitional uncertainty} (which $R$, which $M$). This dominates
    and is far larger: Hubble radius vs.\ particle horizon moves $R$ by tens of
    percent, and baryonic vs.\ matter vs.\ total mass moves $M$ by a factor
    $\sim 20$ (Tables~\ref{tab:R_choices}--\ref{tab:M_choices}). The prefactor
    lands on $4/5$ only for one mutually consistent critical-density choice ---
    which is exactly the Mach--Sciama statement, and exactly why this is not an
    independent measurement of $G$.
\end{itemize}

A genuine prefactor test would require an independent constraint on
$\Muniv/\Runiv$ that does not pass through $G$. We are not aware of one
that is currently tight enough to discriminate, say, $4/5$ from $1/2$ or
$1$ at the few-percent level. The substantive falsifiable content is therefore
the discrete $4/5$-vs-$16/15$ choice (Section~\ref{sec:weight_forced}) and the
$G$-free electron-scale test (Section~\ref{sec:status_predictions}).

\section{Status of the Predictions}
\label{sec:status_predictions}

Honest accounting of what the mechanism currently predicts cleanly,
predicts structurally, and does not yet predict.

\subsection{Cleanly Predicted (G-Free): Bohr Magneton}

Applied at the electron scale, the same helical-jerk mechanism predicts
\begin{equation}
  \mu_B = \frac{e \hbar}{2 m_e}
  \label{eq:bohr_magneton_ref}
\end{equation}
exactly, with no fitted parameters and using only $e$, $\hbar$, $m_e$,
and $c$ --- none of which involve $G$. This is the cleanest test of the
mechanism and is treated in detail in the companion paper.

\subsection{Structurally Predicted (Quantitatively Off): Frame Dragging}

The boost generator structure (BW theorem) implies that rotating masses
drag the local entanglement field, producing a Lense--Thirring-like
effect. The mechanism predicts the correct \emph{structure}: frame
dragging is a boost effect of the modular flow under rotation, not a
separate phenomenon.

The \emph{quantitative} match is currently off by roughly a factor of
$10^3$. A naive collective treatment (random walk of $\sqrt{N_{\mathrm{Earth}}}$
nucleons) does not reproduce $\omega_{\mathrm{FD}}^{\mathrm{GR}} =
2GJ/(c^2 R^3) \approx 1.9 \times 10^{-14}$ rad/s. Closing this gap
requires a more precise treatment of the collective contribution at
the nucleon level, and is open work.

\subsection{Order-of-Magnitude: Dark Energy}

The unresolved entanglement field at the cosmic scale gives an energy
density of order $\rho_{\mathrm{DE}} \approx 8.1 \times 10^{-11}$ J/m$^3$,
compared to the observed $\rho_{\mathrm{DE}}^{\mathrm{obs}} \approx
6.9 \times 10^{-10}$ J/m$^3$. The prediction is within an order of
magnitude. The factor of $\sim 8$ is consistent with the holographic
picture (dark energy stored on the cosmic horizon, a 2D surface,
rather than distributed through the bulk volume), but a precise
holographic accounting is beyond the present scope.

\subsection{Not a Free Prediction: $\alpha_{\mathrm{grav}}$}

The standard gravitational coupling
$\alpha_{\mathrm{grav}} = G m_p^2/(\hbar c) \approx 5.9 \times 10^{-39}$
follows from the predicted relation for $G$ together with $m_p$, $\hbar$,
and $c$. It is therefore consistent with the framework but not an
independent prediction --- it inherits the same circularity as $G$ itself.

\section{Toward a Local Derivation}
\label{sec:local_derivation}

Section~\ref{sec:local_test} showed that the cosmic relation
$G = (4/5)\, c^2 R/M$ is not local. The right next step, if the
mechanism is to be more than a fit at one scale, is to attempt a
local derivation in the spirit of the Jacobson--Padmanabhan--Verlinde
program. We sketch what would be required and what is currently
missing.

\subsection{The Jacobson Program}

\citet{jacobson1995} showed that Einstein's equation can be derived
from the Clausius relation $\delta Q = T\,dS$ applied to local
Rindler horizons, given the Bekenstein--Hawking entropy
$S = A/(4 \ell_P^2)$. \citet{padmanabhan2010} and \citet{verlinde2011}
extended this to ``entropic gravity'' programs in which Newtonian
gravity emerges from local entanglement entropy variation across a
screen.

Crucially, in all of these programs $G$ enters through the
\emph{coefficient} of the area law, $1/(4 \ell_P^2)$. None of them
\emph{derives} the value of $G$; they derive Einstein's equation
\emph{given} the entropy--area coefficient. The Planck length itself
contains $G$ by definition ($\ell_P = \sqrt{\hbar G/c^3}$). So the
deepest known programs of this type still require an external input
that fixes the gravitational coupling.

\subsection{What a Local Version of Our Mechanism Would Need}

For our specific entanglement-torque mechanism to constitute a local
derivation of $G$, three things would have to be true:

\begin{enumerate}
  \item \textbf{A universal entanglement bond density.} There must
    exist a local quantity --- call it $\eta$, with units of
    [energy] $\cdot$ [length]$^{-3}$ $\cdot$ [pair]$^{-1}$, or
    equivalent --- that is the same in every vacuum patch and is
    determined by QFT alone.
  \item \textbf{Newtonian acceleration from torque-of-jerk integrated
    over the patch.} The same orthogonal-jerk $+$ helical-trajectory
    geometry, summed over the local entanglement structure, must
    produce $\mathbf{a} = -\nabla\Phi$ with $\nabla^2 \Phi = 4\pi G\,
    \rho$ in the weak-field limit, with a specific dimensionless
    proportionality constant that does \emph{not} depend on the
    patch's size or contents.
  \item \textbf{$\eta$ computed without $G$.} The crucial step. The
    entanglement bond density $\eta$ must be derivable from purely
    QFT inputs (modular Hamiltonians, area-law coefficients) without
    using $G$ as an input. If $\eta$ requires the Planck length to
    define its cutoff, the program inherits the same circularity as
    Jacobson's.
\end{enumerate}

We can offer (1) and (2) at the level of structural sketch. We do
\emph{not} have (3), and to our knowledge nobody else does either.
This is the open problem that, if solved, would convert the present
program from a consistency check into a derivation.

\subsection{A Concrete Attempt at $\eta$ (and Where It Fails)}

To make the open problem concrete rather than a sketch, here is the most
natural attempt to compute $\eta$ from the bond picture, carried until it
breaks. Take the entropy--area law $S = \eta\, k_B\, \mathcal{A}$ and suppose
the horizon is threaded by entanglement links, each carrying of order one bit
($\sim k_B \ln 2$), with an areal density set by a single UV cutoff length
$\ell$:
\begin{equation}
  \eta \sim \frac{\ln 2}{\ell^2}.
\end{equation}
Feeding this into the Jacobson relation $G = c^3/(4\hbar\eta)$ gives
\begin{equation}
  G \sim \frac{c^3 \ell^2}{4 \hbar \ln 2}.
\end{equation}
Demanding the measured value of $G$ then \emph{fixes the cutoff}:
\begin{equation}
  \ell^2 \sim \frac{4 \hbar G \ln 2}{c^3} = 4\ln 2\;\ell_P^2
  \quad\Longrightarrow\quad
  \ell \sim \ell_P = \sqrt{\frac{\hbar G}{c^3}}.
\end{equation}
This is exactly where the attempt fails. To get the right $G$ the link density
must be the Planck density, $\ell \sim \ell_P$ --- but $\ell_P$ contains $G$ by
definition. We have not computed $\eta$ from $G$-free QFT inputs; we have
re-expressed $G$ in terms of a cutoff that itself encodes $G$. The bond energy
$E_{\mathrm{bond}} = \hbar c/\ell$ at this cutoff is the Planck energy, which
tells the same story. This is the identical circularity that defeats the
Jacobson--Padmanabhan--Verlinde programs at the step of fixing the area-law
coefficient: nothing in the local entanglement counting, as currently
formulated, sets the cutoff without already knowing $G$. A genuine derivation
needs an independent, $G$-free determination of either $\ell$ or $\eta$ --- for
instance from a finite, regulator-independent entanglement entropy density of
the relevant QFT --- which we do not have.

\subsection{Why the Bohr Magneton Is the Real Local Test}

Although a local derivation of $G$ is an open problem, the same
mechanism applied at the electron Compton scale produces an exact,
local, $G$-independent result: $\mu_B = e\hbar/(2 m_e)$ (companion
paper). This is the real evidence that the mechanism is on solid
local ground; what's missing is the bridge from ``the mechanism is
locally real'' to ``the mechanism, applied to the cosmic ensemble,
fixes $G$.''

The honest position is therefore: the local seed is proven (Bohr
magneton). The orchard's shape (Mach--Sciama at cosmic scale) is
consistent with the same seed and with a specific predicted
prefactor, but does not by itself constitute a derivation of $G$.

\section{Computational Verification}

All derivations are implemented in Python scripts:

\begin{itemize}
  \item \texttt{derive\_G.py}---Main $G$ derivation
  \item \texttt{derive\_bond\_energy.py}---$E_{\mathrm{bond}} = \hbar c/r$ derivation
  \item \texttt{derive\_3d\_extension.py}---3+1D extension via Huerta \& van der Velde
  \item \texttt{verify\_symbolic.py}---Symbolic (\texttt{sympy}) proof of the
    $\av{1/r}$ integrals, the discrete $4/5$ vs $16/15$ prefactors, and the
    energy-balance chain (showing $\hbar$, $m_p$, $\kappa$ cancel exactly)
  \item \texttt{gfree\_inputs.py}---Evaluates the prefactor on purely $G$-free
    inputs (expansion-kinematics radii; CMB$+$BBN baryon census), quantifying
    the circularity
  \item \texttt{constants.py}---Physical constants and helper functions
\end{itemize}

All scripts run cleanly and produce consistent results. The symbolic check makes
the core arithmetic exact rather than numerical: run
\texttt{python3 verify\_symbolic.py}.

\bibliographystyle{plainnat}
\bibliography{references}

\end{document}
